\documentclass[10pt,a4paper]{amsart}
\usepackage[margin=19mm]{geometry}
\usepackage{amsmath,amssymb,mathtools}
\usepackage[scr=rsfs,cal=euler]{mathalfa}
\usepackage{xfrac}
\usepackage[unicode,hidelinks,bookmarks=false]{hyperref}
\newcommand{\ie}{i.e.\ }
\newcommand{\cf}{cf.\ }
\newcommand{\resp}{respectively\ }
\newcommand{\Subsection}{\subsection}
\providecommand{\nobreakdash}{\nobreak\hbox{-}}
\setlength{\emergencystretch}{2em}
\multlinegap0pt
\usepackage{bm}
\usepackage{bbm}
\usepackage{dsfont}
\usepackage{xspace}
\usepackage{cancel}
\usepackage{mathtools}
\usepackage{amsthm}
\makeatletter
\@ifundefined{lemma}{\newtheorem{lemma}{Lemma}}{}
\makeatother
\title[Carleson families of cubes related to porous sets]{Carleson families of cubes related to porous sets}
\author{Andrei V. Vasin}
\date{Source: arXiv:2604.06964v1 (CC BY 4.0)}
\begin{document}
\maketitle

\noindent\emph{Source and licence.} Carleson families of cubes related to porous sets. Version arXiv:2604.06964v1 (CC BY 4.0), distributed under the Creative Commons Attribution 4.0 International licence, as recorded in the arXiv licence metadata for this exact version and shown on the abstract page https://arxiv.org/abs/2604.06964v1. Full rights notice: licenses/arxiv-2604.06964-CC-BY-4.0.txt.\par\smallskip

\noindent\textit{Editorial scope.} Counted unit: the Lemma whose source label is \texttt{lem} in the article (source0.tex lines 344--347, source label \texttt{lem}), delivered together with the article's own complete original proof of it (source0.tex lines 348--387, one \texttt{proof} environment of 40 lines). No mathematics is written, repaired or re-derived: statement and proof are the article's own text line for line. Required context retained uncounted: none. Every definition, display and label of those lines is retained unchanged. Omitted from this package and not redistributed: 1--343, 388--828. Nothing inside a retained excerpt is omitted or altered: every retained body is delivered exactly as the article's lines are written, so no omission receipt of any kind accompanies this package. Citations: the retained excerpts carry no citation command at all, so no bibliography entry of the article is transcribed and no reference is redistributed. Reference adaptation: none was needed and none was made; every reference inside the delivered unit resolves inside it, so no unresolved reference or citation is produced. Printed slips kept literal: no printed word, symbol, hypothesis or number of the counted statement or of its proof was altered. Layout and class: the article's own class is \texttt{amsart} with packages including amssymb, amsmath, amsthm, esint, fontenc, inputenc; the wrapper is the lane's pinned \texttt{amsart} editorial wrapper at 10pt on A4, which restates the theorem-like environments the retained text needs and adds the article's own macro definitions that the retained text needs (none), with extra wrapper packages (bm, bbm, dsfont, xspace, cancel, mathtools). The delivered numbering is the unit's own. Assets: no retained span contains a figure environment, a graphics-inclusion command, a PDF-page inclusion, a table or a TikZ picture; no image file is copied into this package, the package declares no asset dependency and no image file is redistributed with it. Licence: version arXiv:2604.06964v1 of this article is distributed under the Creative Commons Attribution 4.0 International licence, as recorded in the arXiv licence metadata for this exact version and shown on the abstract page https://arxiv.org/abs/2604.06964v1. A full rights notice is delivered at licenses/arxiv-2604.06964-CC-BY-4.0.txt. Characters: every retained excerpt is pure ASCII, so the delivered engine, XeLaTeX with its default Latin Modern text font, reports no missing character.
 \begin{lemma}\label{lem}
   Let  $E$ be a porous set with the porosity constant $\eta$. Given  a cube $R\in \mathcal{D}$. Then there is a  disjoint family of free cubes
   \[\mathcal{M}(R) = \{M(Q) \subset Q: Q \in \mathcal{D}_E(R);\;  2^d \eta | M(Q)| \geq |Q|\}.\]
    \end{lemma}

   \begin{proof}
    For $j=0,1,\dots$ let
\[\mathcal{D}_{j,E}(R)=\{Q \in \mathcal{D}_j(R): \:  Q \cap E \neq \varnothing\}.\]
Arguing by induction,  assume that for an integer $ N\geq 0$ there exists a pairwise disjoint family of free cubes
    \[ \mathcal{M}_N(R)=\{M(Q): \;  Q \in \bigcup _{0 \leq n \leq N}\mathcal{D}_{n,E}(R)\},\]
     with the properties
      \begin{enumerate}
         \item[(a)] $ 2^d \eta | M(Q)| \geq |Q|$;
          \item [(b)] Each $Q \in \mathcal{D}_{N,E}(R)$  in addition to    $M(Q)$ may contain no more than one cube $M(Q')$, assigned to  some ancestor $Q' \in \mathcal{D}_{n,E}(R)$, $ 0 \leq n <N$, $Q' \supset Q$.  In this  case there are two distinct children $c_1Q, c_2Q \in  \mathcal{D}_1 (Q)$  such that $c_1Q\supset M(Q)$ and    $c_2Q\supset M(Q')$, respectively.
  \end{enumerate}
 Observe, that the statement is obvious for $N=0$, since we simply take a maximal free cube $M(R)$.

   We check the properties (a)--(b) for $N+1$. For this consider   $\mathcal{D}_1 (Q) \subset\mathcal{D}_{N+1}(R)$ for each $Q \in \mathcal{D}_N(R)$ and choose the assigned free cube for each child in $ \mathcal{D}_{1,E} (Q)$.

   For those children $cQ \in \mathcal{D}_{1,E} (Q)$ such that  $cQ \neq c_1Q$ and  $cQ \neq c_2Q$, we simply choose the largest  free cube $M(cQ) \subset cQ$, $\eta |M(cQ)| \geq  |c Q|$.

   Then, consider the marked children  $c_1Q $ and $c_2Q$, if the latter exists.
    The argument is similar in  both cases.

 Take $c_1Q$ say,  by induction assumption, it already contains free cube
    $c_1Q \supset  M(Q)$.   If    $c_1Q=M(Q)$, that is $c_1Q \notin \mathcal{D}_{1,E} (Q)$, then there is nothing more to be done for $c_1Q$.

  Alternatively,
  if $M(Q) \varsubsetneq c_1Q$,  and recalling that  $M(Q)$ is already assigned to $Q$, we need to find the second  largest free cube $M(c_1Q) \subset c_1Q$  in order to assign it to $c_1Q$.
                     For this  consider  the family $\mathcal{D}_1(c_1Q) \subset \mathcal{D}_{N+2}(R)$. Take a sub-child cube  $ c(c_1Q) \in  \mathcal{D}_{1,E}(c_1Q)$ such that   $c(c_1Q) \cap M(Q) = \varnothing$.
  Choose the largest free cube in  $c(c_1Q)$, but assign it to $c_1Q$,  thus,  denoting it  by $M(c_1Q)$. It holds
     \[\eta |M(c_1Q)| \geq |c(c_1Q)|= 2^{-d}  |c_1Q|\]
    and thus,
     \[2^d \eta |M(c_1Q)| \geq    |c_1Q|\]
     with
 worsening a porosity constant.





     Arguing in the same way for $c_2Q$, we obtain
   all properties (a)--(b) for all cubes from $Q \in \mathcal{D}_{N+1,E}(R)$.
The proof of the lemma is completed.
   \end{proof}

\end{document}
